\documentclass[10pt,a4paper]{amsart}
\usepackage[margin=19mm]{geometry}
\usepackage{amsmath,amssymb,mathtools}
\usepackage[scr=rsfs,cal=euler]{mathalfa}
\usepackage{xfrac}
\usepackage[unicode,hidelinks,bookmarks=false]{hyperref}
\newcommand{\ie}{i.e.\ }
\newcommand{\cf}{cf.\ }
\newcommand{\resp}{respectively\ }
\newcommand{\Subsection}{\subsection}
\providecommand{\nobreakdash}{\nobreak\hbox{-}}
\setlength{\emergencystretch}{2em}
\multlinegap0pt
\usepackage[shortlabels]{enumitem}
\usepackage{amssymb}
\usepackage{bbm}
\usepackage{float}
\usepackage{nicefrac}
\newtheorem{lemma}{Lemma}
\newcommand{\To}{\rightarrow}
\newcommand{\mcd}{\mathcal D}
\title{On the distribution kernels of Toeplitz operators on CR manifolds}
\author{Chin-Yu Hsiao, Ood Shabtai}
\date{Source: arXiv:2512.16506v2 (CC BY 4.0)}
\begin{document}
\maketitle

\noindent\emph{Source and licence.} On the distribution kernels of Toeplitz operators on CR manifolds. Version arXiv:2512.16506v2 (CC BY 4.0), distributed by arXiv under the Creative Commons Attribution 4.0 International licence, http://creativecommons.org/licenses/by/4.0/, as recorded in the arXiv licence metadata for this exact version and shown on the abstract page https://arxiv.org/abs/2512.16506v2. Full licence notice: \texttt{licenses/arxiv-2512.16506-CC-BY-4.0.txt}.\par\smallskip

\noindent\textit{Editorial scope.} Counted unit: the article's lemma l-gue251206yyd (source0.tex lines 1196--1208). The article's complete original proof of it is delivered with it (source0.tex lines 1210--1250, one \texttt{proof} environment of 41 lines). No mathematics is written, repaired or re-derived: the statement and the proof are the article's own lines, in the article's own notation, and no retained line is substituted or re-flowed.  Retained uncounted as the source-local context the proof needs: lines 1157--1164.  Nothing was omitted from any retained span, so the package carries no omission record.  No retained line contains a citation, so the wrapper carries no bibliography and no bibliography entry.  No retained line contains a figure environment, an image-inclusion command or drawing code, so no image is delivered and the package declares no asset dependency.  Editorial formatting only, in addition: an AMS \texttt{amsart} preamble that restates the article's own declarations and macros used by the retained lines, namely the environments \texttt{lemma} on separate counters, together with the standard packages those lines need.  The retained environments print in this document as Lemma 1 in order of appearance, whereas the article prints the same items with the numbering of its own full document.  The complete native source of the article as distributed in the arXiv e-print for this exact version is bundled unchanged as \texttt{sources/191243/source0.tex}.
 \begin{equation}\label{e-gue251206yyd}
			\begin{split}
				&\mbox{$\square_b\hat{A}+\hat{S}=I$ on $U$},\\
                &\mbox{$\hat{A}^*\square_b+\hat{S}^*=I$ on $U$},\\
				&\square_b\hat{S}\equiv0\ \ \mbox{on $U$},\\
				&\hat{S}\equiv \hat{S}^*\equiv \hat{S}^2\ \ \mbox{on $U$},
			\end{split}
		\end{equation}

\begin{lemma}\label{l-gue251206yyd}
Suppose  that \[\square_b: {\rm Dom\,}\square_b\subset L^2(X)\rightarrow L^2(X)\]
has closed range and the partial inverse $N$ is continuous:
\[N: C^\infty(X)\rightarrow C^\infty(X).\]
Let $(U, V, \Gamma, \Psi)$
be an orbifold chart of $X$. Assume that $X$ is strictly pseudoconvex on $U$. Then, for any $\chi\in C^\infty_0(U)$, $\tau\in C^\infty(X)$ with 
${\rm supp\,}\chi\cap{\rm supp\,}\tau=\emptyset$, we have 
\begin{equation}\label{e-gue251209yyds}\begin{split}
    &\chi\Pi\tau\equiv0\ \ \mbox{on $X$},\\
    &\tau\Pi\chi\equiv0\ \ \mbox{on $X$}.
\end{split}
\end{equation}
\end{lemma} 

\begin{proof}
Let $\hat A$ and $\hat S$ be properly supported continuous operators:
    \[\hat A, \hat S: C^\infty(U)\To C^\infty(U)\]
    given by \eqref{e-gue251206yyd}. We have
    \begin{equation}\label{e-gue251209yyd}
			\begin{split}
	&\mbox{$\square_b\hat{A}+\hat{S}=I$ on $U$},\\
    &\mbox{$\hat{A}^*\square_b+\hat{S}^*=I$ on $U$},\\
				&\square_b\hat{S}\equiv0\ \ \mbox{on $U$},\\
				&\hat{S}\equiv \hat{S}^*\equiv \hat{S}^2\ \ \mbox{on $U$},
			\end{split}
		\end{equation}
		where $\hat{S}^*$ and $\hat{A}^*$ are the formal adjoints of $\hat{S}$ and $\hat{A}$
		with respect to $(\,\cdot\,|\,\cdot\,)$ respectively.  
        Let $\chi\in C^\infty_0(U)$. 
From the second equation of \eqref{e-gue251209yyd} and notice that $\hat A$ and $\hat S$ are properly supported , we have 
\begin{equation}\label{e-gue251209yydI} \mbox{$\chi\hat{S}^*\Pi=\chi\Pi$ on $X$}
\end{equation}
and hence
\begin{equation}\label{e-gue251209yydII}
\mbox{$\Pi\hat{S}\chi=\Pi\chi$ on $X$}.\end{equation}

On the other hand, we have 
\begin{equation}\label{e-gue251209yydIII}\hat S\chi=N\square_b\hat S\chi+\Pi\hat S\chi.\end{equation}
Since $N$ maps smooth functions to smooth functions,
\[N\square_b\hat S\chi: \mcd'(X)\To C^\infty(X)\]
is continuous. Hence 
\[N\square_b\hat S\chi\equiv0.\]
From this observation and \eqref{e-gue251209yydIII}, we get 
\begin{equation}\label{e-gue251209yyda}\hat S\chi=\Pi\hat S\chi+F\ \ \mbox{on $X$},\end{equation}
where $F\equiv0$ on $X$. 

From \eqref{e-gue251209yydII} and \eqref{e-gue251209yyda}, we get 
\[\hat S\chi=\Pi\chi+F\ \ \mbox{on $X$}\]
and hence 
\begin{equation}\label{e-gue251209yydb}
\tau\hat S\chi=\tau\Pi\chi+\tau F,
\end{equation}
where $\tau\in C^\infty(X)$ with 
${\rm supp\,}\chi\cap{\rm supp\,}\tau=\emptyset$. From \eqref{e-gue251209yydb} and notice that $\hat S$ is smoothing away the diagonal, the lemma follows. 
\end{proof}

\end{document}
